\scanpage{312}

\mspage{51}
% archive sheet number: 311

de $\underline{\mathrm{Ty}}$ :
\[
f_\tau:I_\tau\unsure{\text{ --- }|\longrightarrow}e,\qquad(\tau\in\mathrm{Ty})
\]
et on a
\[
f_\tau*f_{\tau'}=f_{\tau\tau'},
\]
i.e. \deleted{\unsure{la multip.}} l'application
\diagram{\begin{tikzcd}[ampersand replacement=\&,row sep=small,column sep=small]
\tau \arrow[r,mapsto] \& f_\tau : \& \operatorname{Ob}(\underline{\mathrm{Ty}}) \arrow[r] \arrow[d,equal] \& \operatorname{Fl}(\underline{\mathrm{Ty}}) \arrow[dd,equal]\\
\& \& \mathrm{Ty} \arrow[d,equal] \& \deleted{\unsure{\mathrm{Ty}}}\\
\& \& \operatorname{Nerf}_0(\underline{\mathrm{Ty}}) \& \operatorname{Nerf}_1(\underline{\mathrm{Ty}})
\end{tikzcd}}
est un homomorphisme de monoïdes, transformant $\Delta_\varepsilon$ $(\varepsilon\in\{+1,-1\})$ dans $\deleted{\unsure{\varepsilon}} f_{\varepsilon,0}$. \deleted{Elle n'a pas de \unsure{mérite|intérêt}, puisque c'est l'application de} \deleted{L'application \unsure{de|se}} Ce n'est pas l'application de dégénérescence de $\operatorname{Nerf}_0$ dans $\operatorname{Nerf}_1$, i.e. l'application $\tau\unsure{\rightarrow|\text{ --- }}\mathrm{id}_\tau$. Celle ci est donnée par
\[
\begin{aligned}
 \operatorname{Ob}(\underline{\mathrm{Ty}})&\longrightarrow \operatorname{Fl}(\mathrm{Ty})\\
 \tau&\mathrel{\mapstochar\relbar}\mathrm{id}_\tau
\end{aligned}
:\quad
\left|\begin{aligned}
\Delta_\varepsilon&\longmapsto f_{\varepsilon,1}\qquad(\varepsilon\in\{+1,-1\})\\
\text{i.e.}\qquad \Delta_+&\mathrel{\mapstochar\relbar} f_{+,1}\\
\Delta_-&\mathrel{\mapstochar\relbar} f_{-,1}.
\end{aligned}\right.
\]

Quant aux homomorphismes source et but
\diagram{\begin{tikzcd}[ampersand replacement=\&,column sep=large]
\operatorname{Fl}(\underline{\mathrm{Ty}})\arrow[r,shift left,"\sigma\ (\text{source})"]\arrow[r,shift right,"\beta\ (\text{but})"']\&\operatorname{Ob}(\underline{\mathrm{Ty}})
\end{tikzcd},}
elles sont \unsure{évidentes} sur les générateurs
\[
\left\{\begin{aligned}
\sigma(f_{\varepsilon,\ell})&=\Delta_\varepsilon\\
\beta(f_{\varepsilon,\ell})&=\Delta_\varepsilon^\ell.
\end{aligned}\right.
\]

Je vais décrire maintenant $\operatorname{Nerf}_2(\mathrm{Ty})$, i.e. les \deleted{\unsure{systè}} systèmes
\diagram{\begin{tikzcd}[ampersand replacement=\&,column sep=normal]
\tau\arrow[r,"f"] \& \tau'\arrow[r,"g"] \& \tau''.
\end{tikzcd}}
